\chapter{Concepts}\label{ch:concepts}

The features in this manual rest on a handful of ideas. Once these are clear, the rest follows: why a
value can be drawn but not edited, why \keys{Ctrl+Alt+B} sometimes pauses and sometimes stops, why a
Grasshopper component needs its script editor to be paused. Each idea gets a short section here; the
chapters after it give the details.

\section{Debugging in two minutes}\label{sec:debugging}

A \emph{debugger} runs your script and stops it at a line you choose. While it is stopped, or
\emph{paused}, you can look at every variable and see what the script really computed, not what you
think it computed.

\begin{xltabular}{\linewidth}{@{}P{0.3\linewidth}L@{}}
\toprule
\textsf{\textbf{Term}} & \textsf{\textbf{What it means}} \\ \midrule
\ui{Breakpoint}            & A red dot left of a line number. Click there to set or remove one. The script pauses when it reaches that line. \\
\ui{Debug} \key{F5}        & Runs the script in the debugger, so it stops at breakpoints. (\key{F5} without a breakpoint is a plain run; miniVivisection's \ui{Debug} button always debugs, \see{sec:runbuttons}.) \\
\ui{Paused}                & The script waits. The highlighted line is the \emph{next} one to run; it has \emph{not} run yet. \\
\ui{Continue} \key{F5}     & Runs on to the next breakpoint, or to the end. \\
\ui{Step over} \key{F10}   & Runs one line, then pauses again. \\
\ui{Step into} \key{F11}   & If the line calls a function, goes into it and pauses at its first line. \\
\ui{Step out} \keys{Shift+F11} & Finishes the current function and pauses where it was called. \\
\ui{Stop} \keys{Shift+F5}  & Ends the run. \\
\ui{Variables} tab         & The editor's list of the paused script's variables, at the bottom. \\
\bottomrule
\end{xltabular}

The method is always the same: guess where things go wrong, set a breakpoint just before that place,
look at the values, step forward, and narrow it down until you see the line that makes a wrong value.
miniVivisection makes the \emph{look} part faster: you see values as geometry, in the code, and in
Rhino at once.

\section{The paused frame}\label{sec:frame}

When the debugger pauses, the script is in the middle of a function, a \emph{frame}, with its local
variables, the globals of its module and everything it imported. Nearly everything miniVivisection
does happens in this frame:

\begin{itemize}
  \item A watch such as \code{curves[0].PointAtEnd} is a Python expression, evaluated in the paused
        frame at every pause. It sees the same names the paused line sees.
  \item A tooltip reads a name \emph{as Python sees it at the hovered line}. Over a global \code{crv}
        at module level you get the global curve, even while paused inside a function with a local
        \code{crv}.
  \item \vn{Python} runs what you type in the paused frame. An assignment changes the script's
        variable, and the script continues with the new value.
  \item Editing in Rhino writes the edited geometry back with an assignment in the paused frame.
\end{itemize}

Because watches run again at every pause, an expression with a side effect, such as
\code{curves.pop()}, would act each time. Keep watches to expressions that only read.

Python evaluation needs the \emph{paused thread}: the frame is found by walking up that thread's
stack. When the script runs on Rhino's UI thread, the plug-in is already there. When it runs on a
thread of its own, the plug-in hands the work to that thread through FreePause's wait loop
(\see{sec:freepause}).

\section{Drawn, not added}\label{sec:drawn}

What \vn{Watch} shows in the viewports is \emph{drawn}: a display conduit paints copies of the values
on top of the scene at every redraw. Nothing is added to the document, nothing is saved in the
\file{.3dm}, and nothing has to be cleaned up. When the run ends, the drawing goes away.

Two features make real Rhino objects instead, on purpose:

\begin{description}
  \item[Bake] copies drawn values into the document, on the layer \file{Script Editor Watch::Baked}, to
        keep them (\see{sec:bake}).
  \item[Edit in Rhino] turns a value into objects on the layer \file{Script Editor Watch::Edit}, so you
        can edit them with Rhino's tools. They are removed when the run ends (\see{sec:editinrhino}).
\end{description}

The same principle applies in Grasshopper: \vn{Edit}'s edit objects belong to the component, not to the
model. They are never saved in the \file{.3dm} and don't mark the model as changed (\see{sec:vedit}).

\section{Write-back}\label{sec:writeback}

Editing in Rhino is a round trip:

\begin{center}
\begin{tikzpicture}[node distance=7mm and 9mm, every node/.style={font=\sffamily\small},
    box/.style={draw=mvbluemid, fill=mvbluelight, rounded corners=3pt, minimum height=9mm, inner xsep=6pt, align=center},
    arr/.style={-{Stealth[length=2.5mm]}, mvblue, line width=0.8pt}]
  \node[box] (var) {Variable in the\\paused frame};
  \node[box, right=22mm of var] (obj) {Rhino object\\(layer …::Edit)};
  \node[box, right=12mm of obj] (edit) {You edit it\\Gumball, points, Move};
  \node[box, below=9mm of obj] (back) {Converted to the variable's type,\\assigned in the paused frame};
  \draw[arr] (var) -- node[above, font=\sffamily\scriptsize, text=mvblue]{Edit in Rhino} (obj);
  \draw[arr] (obj) -- (edit);
  \draw[arr] (edit) |- node[near end, below, font=\sffamily\scriptsize, text=mvblue]{a command ends} (back);
  \draw[arr] (back) -| node[near end, left, font=\sffamily\scriptsize, text=mvblue]{the script continues with it} (var);
\end{tikzpicture}
\end{center}

The edited geometry is turned back into the variable's own type: a moved point into a \code{Point3d}, a
curve that is still a line into a \code{Line}, and so on. A list item is written back as
\code{curves[2] = …}. Rhino's Undo and Redo write the earlier or later geometry back too.

\begin{caution}[A write-back lasts for this run]
What you change while the script is paused, by dragging points or typing into \vn{Python}, changes
the running script's data, not its code. When you have found the fix, change the code.
\end{caution}

\section{Rhino's UI thread and FreePause}\label{sec:uithread}

Rhino has one \emph{UI thread}: it draws the viewports, reads the keyboard and mouse, and runs
commands. A script started from the Script Editor normally runs on that same thread. When such a
script pauses, Rhino's own debugger keeps the thread busy in a \ui{Select Action} prompt on the command
line: typed text goes to that prompt, objects can't be selected, and only toolbar macros get through.

\emph{FreePause} replaces that prompt. While the script is paused, Rhino sits at its normal
\code{Command:} prompt and keeps its message loop going. You can select, run commands, use the Gumball
and orbit the view. You continue and step with the Script Editor's buttons or keys
(\key{F5}, \key{F10}, \key{F11}, \keys{Shift+F11}, \keys{Shift+F5}).

A script can also run on a thread of its own: add \code{\# async:true} after its first line
(\code{\#! python 3}). Rhino then stays free even while the script runs, not only while it is paused.
Such a script mustn't use Rhino's UI directly. The difference matters for threads (\see{sec:freezeguard})
and for Grasshopper (\see{sec:ghasync}).

\section{Debug runs and plain runs: pause or stop}\label{sec:runkinds}

A script can run in two ways, and the difference decides what \keys{Ctrl+Alt+B} can do with it.

\begin{description}
  \item[A debug run] is traced by the debugger line by line. It stops at breakpoints, and
        \keys{Ctrl+Alt+B} \emph{pauses} it at its next line: you can look at it, step, and continue.
  \item[A plain run] runs at full speed, without the debugger. Breakpoints are ignored. It can't be
        paused, but with miniVivisection \keys{Ctrl+Alt+B} can \emph{stop} it, and the plug-in records
        where it was: the file, the line, the call stack and the values of that moment. The Script
        Editor later opens the script at that line and explains what happened (\emph{Stop and record}).
\end{description}

\begin{center}
\begin{tikzpicture}[every node/.style={font=\sffamily\small},
    run/.style={draw=mvbluemid, fill=mvbluelight, rounded corners=3pt, minimum width=36mm, minimum height=9mm, align=center},
    res/.style={draw=#1, fill=#1!8, rounded corners=3pt, minimum width=52mm, minimum height=9mm, align=center},
    arr/.style={-{Stealth[length=2.5mm]}, line width=0.8pt}]
  \node[run] (dbg) at (0,0) {\textbf{Debug run}\\Debug button, F5 with a breakpoint};
  \node[run] (pln) at (0,-2.1) {\textbf{Plain run}\\Run, RunPythonScript, Grasshopper};
  \node[res=mvgreen] (pause) at (7.2,0) {\textbf{pauses} at the next line\\look, step, continue};
  \node[res=mvorange] (stop) at (7.2,-2.1) {\textbf{stops}, and records where\\the editor opens at that line};
  \draw[arr, mvgreen] (dbg) -- node[above]{Ctrl+Alt+B} (pause);
  \draw[arr, mvorange] (pln) -- node[above]{Ctrl+Alt+B} (stop);
\end{tikzpicture}
\end{center}

Which kind a script gets depends on how it is started:

\begin{itemize}
  \item In the Script Editor, \key{F5} debugs only when the script has a breakpoint, otherwise it is a
        plain run. miniVivisection's \ui{Run} and \ui{Debug} buttons make it explicit (\see{sec:runbuttons}).
  \item \code{RunPythonScript} and \code{ScriptEditor Run} run plainly, unless their \emph{Debug} mode is chosen
        (\see{sec:runpythonscript}).
  \item Grasshopper script components run plainly, unless \emph{Always debug} is on or the run was
        started from the component's script editor with \ui{Debug} (\see{sec:ghmodes}).
\end{itemize}

\keys{Ctrl+Alt+B} is a Windows hotkey, taken only while a script runs. A script on Rhino's UI thread
holds that thread, so nothing on screen can change; the hotkey answers with a sound instead. A soft
sound means it paused or stopped something; an error sound means there was nothing it could do.

\section{Threads}\label{sec:threadsconcept}

Python scripts can start threads (\code{threading.Thread}). The debugger traces each of them. When one
thread reaches a breakpoint and pauses, the others stop at their next line, until you continue.
\vn{Threads} lists every thread and where it is.

A paused \emph{worker} thread is dangerous when the script runs on Rhino's UI thread: the UI thread
waits for the worker (in \code{join()}, say), the worker waits for you, and you can't click anything,
because the UI thread is the one that would handle the click. Rhino freezes. miniVivisection's
\emph{freeze guard} skips such a pause and says why (\see{sec:freezeguard}). To debug threads, run
the script with \code{\# async:true}.

\section{Grasshopper: components run inside a solution}\label{sec:ghconcept}

A Grasshopper script component is not a script you run once. It runs once per item or per branch of
its input data, inside a Grasshopper \emph{solution}, on Rhino's UI thread. Three things follow:

\begin{description}
  \item[Runs have names.] A component that fails for some data fails in particular runs. miniVivisection
        names them by run number, branch (\code{\{0;3\}}) and item (\code{\{0;3\}[2]}), and a breakpoint can
        pause for exactly that data (\see{sec:failingruns}).
  \item[The debugger lives in the component's editor.] A component's breakpoints and stepping belong to
        its script editor window. With the editor closed, there is no debugger to pause in, so
        \keys{Ctrl+Alt+B} can only stop the component. Under \emph{Always debug}, miniVivisection gives
        every Python~3 component a hidden editor, so it can pause even if you never opened it
        (\see{sec:hiddeneditors}).
  \item[A pause holds the solution.] While a component is paused, its solution is unfinished.
        Grasshopper doesn't allow recomputing during a solution. miniVivisection keeps the canvas usable,
        computes changes downstream with the paused component's last data, marks what depends on the
        paused component, and runs everything properly when the pause ends (\cref{ch:ghpaused}).
\end{description}

The three modes for script components choose between these behaviors: \emph{Normal} (plain runs,
left alone), \emph{Stop and record} (plain runs, \keys{Ctrl+Alt+B} stops a hang, the default) and
\emph{Always debug} (every recompute in the debugger).

\section{Markdown as documentation}\label{sec:markdownconcept}

Markdown is plain text with light marks: \code{\# Heading}, \code{**bold**}, \code{*italic*},
\code{`code`}, \code{- list}, \code{[text](target)}. miniVivisection uses it in three places, with the same
parser and the same link targets:

\begin{itemize}
  \item \textbf{Panels in Grasshopper} show Markdown rendered and scripts as code (\cref{ch:panels}).
  \item \textbf{Markdown text in Rhino}, made by the \code{MarkdownText} command, is real Rhino text that
        you edit by double-clicking it (\cref{ch:markdowntext}).
  \item \vn{Text} turns a Panel's text into Rhino text in Grasshopper (\see{sec:vtext}).
\end{itemize}

Links connect these documents with what they describe: a note on the canvas can lead to a
component, a Rhino object, a named view, another definition, a model or a script, and back
(\cref{ch:links}). The lessons use this: each ends with links to the one before and after, and an
overview links to all of them.

\section{Prompts and AI}\label{sec:aiconcept}

Three prompts work on the paused script, in the Script Editor and in a window of their own for
Grasshopper:

\begin{description}
  \item[\vn{Python}:] Python in the paused frame. Try a fix on live data before you change the code.
  \item[\vn{PowerShell}:] A shell in the script's folder, for files, git and tools.
  \item[\vn{Claude}:] A chat with Claude Code. Claude can read the script, the paused values and the run's
        errors, set breakpoints, step, propose edits and, with RhinoAI, work in the model and in
        Grasshopper.
\end{description}

Claude acts only through tools the plug-in gives it, and you keep control. Code it wants to run waits
for your \code{/run}, edits to a paused script wait until the run ends, RhinoAI's destructive tools stay
off until \code{/allow on}, and in Grasshopper each answer is one undo step (\see{sec:claude}).

\section{The V names and the two switches}\label{sec:vnames}

Everything miniVivisection adds is named with a V, so it is told apart from the editor's own:
\vn{Watch}, \vn{Threads}, \vn{Python}, \vn{PowerShell}, \vn{Claude}, \vn{Log}, \vn{Settings}, \vn{Help},
the Grasshopper components \vn{Edit}, \vn{Text} and \vn{Protect}, and the switches \vn{AI} and \vn{Debug}.
Its items in Grasshopper's menus start with \ui{V°} too (\ui{V° Show as}, \ui{V° Debug this component}).

The mark between the V and the name is a setting (\setting{Mark in the V names}): \ui{°} (the default),
\ui{|}, \ui{∣} (a bar that stays on the line), or either bar with thin spaces (\ui{V\,∣\,Watch}). It applies
to Script Editors opened afterwards, and to \vn{Edit} and \vn{Text} components placed afterwards. This manual
writes V°Watch. To search for it, type V° (\keys{Alt+0176} on the number pad) or just the name.

Two switches in the Script Editor's top bar turn miniVivisection on and off as a whole:

\begin{xltabular}{\linewidth}{@{}P{0.18\linewidth}L@{}}
\toprule
\vn{Debug} & Every debugging feature: the V tabs and pages, drawing, tooltip values, line-end values,
             the hover preview, watch expressions, conditional breakpoints (off: every breakpoint pauses),
             FreePause, the menu items. Off, they are all gone; on, everything is back as it was, with the
             settings unchanged. \\
\vn{AI}    & The \vn{Claude} tab and page. Off, they are hidden in every Script Editor; their sessions keep running. \\
\bottomrule
\end{xltabular}
